% PAPER_MATTER_GRAVITY --- standalone Physical Review D article.
%
% Focused "bullet" paper in the style of PAPER_GOLDSTONE_PRD and
% PAPER_BTFR_MNRAS: small, self-contained, one isolable claim. The claim is the
% chain that survives independently of the (retracted) emergent-photon question:
% a pure causal substrate with a single internal orientation field produces
%   (i)   an emergent Newtonian gravitational sector (theta ~ 1/r, Poisson,
%         linear superposition, Schwarzschild dilation by chain counting),
%   (ii)  a stable topological particle (the SU(2) Skyrmion, B=1, spin-1/2,
%         fermionic) that is QUANTITATIVELY a baryon (one calibration fixes e;
%         mu_p/mu_n, radius, multiplet degeneracies parameter-free vs ANW), and
%   (iii) whose OWN energy-density profile, fed as the literal source into the
%         network's gravitational relaxer, sources that very field (theta ~ M/r)
%         --- the heart of the paper.
% Sections II (substrate/protocol) and the two theorems (group elimination;
% Cauchy-Schwarz Skyrme bound) are reused, with adjustments, from the approved
% text of PAPER_GOLDSTONE_PRD / TEIC Document I.
%
% UPDATE (jun/2026): incorporates two campaigns the first draft named as future
% work. MG1 ("Skyrmion as a direct source of gravity"): the soliton's own
% relaxed eps(r) profile is fed as the literal source, closing the Sec.-VII gap
% from a COMPOSITION of two measured laws to a SINGLE direct measurement
% (exponent -0.992, G_net constant to 5 digits, = top-hat to 0.02%). BQ
% ("quantitative baryon phenomenology"): collective-coordinate quantization with
% the full (sigma+Skyrme) inertia; one N-Delta calibration -> e=5.39 (1% of ANW),
% then mu_p/mu_n=-1.51, radius 0.65 fm, degeneracies [4,16,36] parameter-free.
%
% Scope is strictly SU(2) + gravity + matter. SU(3)/confinement, MOND, and dark
% matter are explicitly out of scope (companion papers). The emergent photon is
% NEVER claimed as a positive result; the wave excitation of the same field is
% mentioned only by reference to the companion Goldstone paper (a scalar, not a
% photon).
%
% Reframe (jul/2026, PLANO_REENQUADRAMENTO): intro motivates the CST matter gap
% (no derived particle/statistics/source-field loop); Sec. VII BTFR paragraph
% reworded as a SEPARATE phenomenological line, not the empirical content of this
% sector. No result or claim changed.
%
% Compiles with pdflatex + bibtex (revtex4-2 + apsrev4-2.bst) against the shared
% TEIC_CONSOLIDATED.bib.
\documentclass[aps,prd,reprint,nofootinbib,superscriptaddress,floatfix]{revtex4-2}
\usepackage{amsmath,amssymb}
\usepackage{amsthm}
\usepackage{bm}
\usepackage{booktabs}
\usepackage{hyperref}

\theoremstyle{plain}
\newtheorem{theorem}{Theorem}
\newtheorem{proposition}{Proposition}

\newcommand{\Sderiv}{\textsf{[Derived]}}
\newcommand{\Sident}{\textsf{[Identified]}}
\newcommand{\Sext}{\textsf{[External]}}
\newcommand{\su}[1]{SU(#1)}

\begin{document}

\title{Topological matter and emergent gravitation from a causal network:\\
a soliton that sources its own gravitational field}

\author{Miqueias Alves Mendes}
\email{mikalvesm@gmail.com}
\affiliation{Independent Researcher, Ibiapina, Cear\'a, Brazil}
\date{\today}

\begin{abstract}
We ask whether a discrete causal substrate carrying a single internal orientation
field can produce, simultaneously and from the same construction, stable matter with
the correct quantum statistics \emph{and} the gravitational field that matter should
generate. Working under a guard that forbids any relativistic expression in the
data-generating code, we establish four measured links. (i)~The coarse-grained
causal-density field relaxes, under the Benincasa--Dowker action with a localised
source and no imposed profile, to $\theta\propto r^{-1.02}$ (exterior exponent
$-1.00$ for every source shape tested, matching an independent Gauss-law quadrature
to $0.9998$); superposition is exact for the linear action (residual $10^{-9}$), and
a non-circular chain-counting measurement of the Schwarzschild dilation reproduces
$\sqrt{1-2M/r}$ to $0.2\%$. (ii)~Among the candidate internal groups, only $\su2$
admits a conserved, localised, finite-energy topological charge ($\pi_3(S^3)=
\mathbb{Z}$); $U(1)$ is excluded by $\pi_2(S^1)=0$, by theorem and by measurement.
(iii)~A pointwise Cauchy--Schwarz bound $K\le\tfrac23 S$ (verified against $10^6$
adversarial configurations) proves the Skyrme stabiliser cannot dominate
spontaneously and is therefore an external input, declared as such. (iv)~The $\su2$
soliton exists with $B=1$, virial mass $M\approx146$--$207$ in lattice units, and
measured spin-$\tfrac12$ fermionic statistics (rigid-rotor band, Finkelstein--
Rubinstein constraint, exchange phase $=2\pi$-rotation phase $\in\mathbb{Z}_2$).
(v)~Quantising the soliton's collective rotation with the full ($\sigma$+Skyrme)
inertia, a \emph{single} calibration---the $N$--$\Delta$ splitting---fixes the one
dimensionless coupling to $e=5.39$ ($1\%$ from Adkins--Nappi--Witten), after which
the magnetic-moment ratio $\mu_p/\mu_n=-1.51$ (experiment $-1.46$), the isoscalar
radius $0.65$~fm and the multiplet degeneracies $(2j+1)^2=[4,16,36]$ are
parameter-free: the soliton is quantitatively a baryon, only the absolute GeV scale
($f_\pi$) remaining external. The central result unites matter and gravity by a
direct measurement: feeding the soliton's \emph{own} relaxed energy-density profile
$\varepsilon(r)$ as the literal source into the same Benincasa--Dowker relaxer yields
$\theta=G_{\rm net}\,M/r$ with exterior exponent $-0.992$ and $G_{\rm net}=A/M$
constant to five digits across the soliton mass, identical to a top-hat of equal mass
to $0.02\%$. The form is derived; the absolute coupling $G_{\rm net}\propto1/K$ and the
GeV scale are external. Matter and gravity are one measurement, not two coincident ones.
\end{abstract}

\maketitle

\section{Introduction}\label{sec:intro}

A recurring ambition of emergent-spacetime programmes is that gravitation and matter
should not be postulated separately but should arise together from one microscopic
structure. We pose a sharp, measurable version of the question for the discrete
substrate of causal set theory (CST): can a causal network equipped with a single
internal orientation degree of freedom produce, from the \emph{same} construction,
(a)~a gravitational sector in which a localised source generates a Newtonian
$1/r$ potential, and (b)~a stable point-like particle with the correct
spin-statistics---and, decisively, does that particle's mass source the field
of~(a)?

In CST, spacetime is fundamentally a locally finite partial order, the order being
the causal relation, and a Lorentzian manifold is an approximation to it
\cite{BLMS1987}. The substrate's defining virtue is that a Poisson sprinkling breaks
no Lorentz frame on average \cite{BombelliHensonSorkin2009}, and from it CST recovers
spacetime dimension and a scalar-curvature (d'Alembertian) operator
\cite{Myrheim1978,BenincasaDowker2010}. We take these geometric results as
established and claim no novelty for them. What CST has so far lacked is the other
half of physics: free scalar fields on a causal set are well developed
\cite{Surya2019}, but no stable particle, no quantum statistics, and no closed
source--field loop have been derived from the causal order itself. That is the gap
this paper addresses: we add one internal field and measure what it forces.

Our results are stated as what the measurements support, and each is tagged with its
epistemic status using the vocabulary of the companion documents: \Sderiv{} (derived
by measurement under the guard), \Sident{} (identified with a known object by
matching quantum numbers, not by an absolute-scale calculation), and \Sext{}
(an external input, declared rather than hidden). We establish an emergent Newtonian
gravitational sector (\S\ref{sec:gravity}), show by elimination that stable
point-like matter forces the internal group to $\su2$ (\S\ref{sec:group}), prove
that the soliton's stabiliser is necessarily external (\S\ref{sec:skyrme}), exhibit
the soliton and measure its mass and fermionic spin-$\tfrac12$ statistics
(\S\ref{sec:soliton}), quantise its collective rotation to reproduce the dimensionless
baryon phenomenology (\S\ref{sec:baryon}), and then close the loop: the soliton's own
energy-density profile sources the gravitational field in proportion to its mass
(\S\ref{sec:sourcesgravity}). We delimit precisely
what is and is not established (\S\ref{sec:claim}) and conclude
(\S\ref{sec:conclusions}). The orientation field belongs to a broader programme
\cite{TEICDoc1}; the wave excitation of the \emph{same} field is the subject of a
separate measurement and is a pair of relativistic internal scalars
\cite{GoldstonePRD}, not addressed here. This paper is self-contained and limited to
the $\su2$ matter-and-gravity sector; $\su3$ confinement, galactic dynamics, and
dark matter are out of scope and treated in separate companion work.

\section{Substrate and protocol}\label{sec:protocol}

\subsection{The causal set}
The substrate is a causal set \cite{BLMS1987}: a finite set of events with a
transitive, irreflexive, locally finite partial order $\prec$, realised by Poisson
sprinkling into $(3{+}1)$-dimensional Minkowski space at intensity $\rho$, with
$x\prec y$ when $y$ is in the causal future of $x$. The \emph{links} are the
relations of the transitive reduction (the Hasse diagram); graph distance along
links is the separation used below. These are the standard CST constructions, used
without modification.

\subsection{Anti-circularity}\label{sec:guard}
Because we claim that relativistic and gravitational structure \emph{emerges}, we
must exclude having inserted it. A guard tokenises every data-generating module and
fails if a Lorentz factor $1/\sqrt{1-\beta^2}$, a gravitational redshift
$\sqrt{1-2M/r}$, or any complex literal appears outside an explicitly labelled,
non-generating comparison block. The speed of light and Newton's constant never
appear in a generator: in the gravitational sector (\S\ref{sec:gravity}) the source
couples through a free, dimensionless weight, never through $G$; in the matter sector
(\S\S\ref{sec:soliton}) the $\su2$ field is carried by unit quaternions---a real
division algebra ($S^3$)---so no complex literal enters and the guard passes
verbatim. Every number below is produced by code that passes this guard; the closed
relativistic forms appear only in quarantined comparison-only blocks, never in a
generator.

\subsection{The internal field}\label{sec:field}
On each event we place an internal orientation and couple causal neighbours by the
minimal rotation-invariant misalignment tension
\begin{equation}
S=\sum_{\langle ij\rangle}\Delta\tau_{ij}\,\big[\,1-\cos\angle(\vec n_i,\vec n_j)\,\big],
\label{eq:sigma}
\end{equation}
with $\Delta\tau_{ij}$ the link proper time. For an $S^2$ orientation this is the
lattice $O(3)$ nonlinear sigma model on the causal Hasse graph; the choice of target
is provisional, and \S\ref{sec:group} shows which group stable matter forces. The
internal symmetry of Eq.~\eqref{eq:sigma} is global and acts only on the internal
index, decoupled from the spacetime causal order.

\section{Emergent Newtonian gravitation}\label{sec:gravity}

We first establish the gravitational sector as a property of the causal-density
field, independently of the matter content; \S\ref{sec:sourcesgravity} will then feed
the soliton into it as a source.

\subsection{Free Monte Carlo: the $1/r$ profile is an output}
\emph{Method.} Write $\theta_i=\delta\rho(r_i)/\rho_0$ for the coarse-grained
causal-density perturbation. We source a localised core and relax $\theta$ freely
under the static Benincasa--Dowker (BD) gradient action with event-number
conservation as an exact constraint,
\begin{equation}
\begin{split}
E[\theta]={}&\frac{K}{2}\sum_i\Big(\frac{\theta_{i+1}-\theta_i}{\Delta r_i}\Big)^2 V_i
\;-\;\kappa\,M\sum_i s_i\,\theta_i,\\
&\sum_i\theta_i V_i=0,
\end{split}
\label{eq:BDenergy}
\end{equation}
with $V_i$ the shell volume, $s_i$ a finite source-core profile, $\kappa$ a free
coupling, and $K$ the action stiffness. \emph{No} functional form is imposed and
neither $G$ nor any relativistic expression enters the generator. \emph{Result.} The
relaxed field is an over-density decaying as $\theta=A/r+C$; the log--log tail
exponent, after removing the exactly predicted conservation offset $C$
($C_{\rm pred}=-0.0257$ vs.\ $C_{\rm fit}=-0.0281$), is $-1.018\approx-1$, the
Newtonian $1/r$. Calibrating $\kappa$ once so that $A=M$ in units $G=c=1$, the
emergent $\rho_{\rm eff}/\rho_0=1+M/r$ reproduces the weak-field
$\rho_0(1+GM/rc^2)$ to a maximum relative error of $0.62\%$ (correlation $0.9991$).
\Sderiv{} for the functional form $1/r$ and the linear law.

\subsection{The law is Poisson, and linear}
\emph{Generic source shapes.} The $1/r$ falloff is not special to a point mass.
Relaxing four extended sources (homogeneous sphere, thin disk, NFW, exponential), the
equilibrium field equals the independent radial-quadrature Poisson solution
$\theta'(r)=-Q_{\rm net}(r)/(4\pi K r^2)$ to correlation $0.9998$, with an exterior
exponent $-0.996$ to $-0.998$ for every bounded shape---the harmonic-exterior
signature of $\nabla^2\theta=J$, shape-independent. \emph{Superposition.} Because the
implemented BD action is quadratic, its Euler--Lagrange equation is the linear
Poisson equation; placing two cores a distance $d$ apart, the noise-free solver gives
$\theta_{\rm tot}=\theta_1+\theta_2$ to a residual $\mathrm{rms}/\sigma=1.8\times
10^{-9}$ (correlation $1.000000$). Superposition is exact: the network obeys a linear
equation, not a nonlinear one in disguise. \Sderiv{}.

\subsection{Schwarzschild dilation by chain counting}
\emph{Method.} Sprinkling the Schwarzschild slice in tortoise coordinates, where the
metric enters the generator only as the proper-volume weight $\Omega^2=(1-2M/r)$
(never square-rooted), we measure the static-observer clock rate $d\tau/dt$ purely by
counting (the estimator $\tau=\sqrt{2N/\rho_0}$). \emph{Result.} The counted clock
rate reproduces $\sqrt{1-2M/r}$ with correlation $0.99999$ and maximum relative error
$0.2\%$ across $3\le r\le160$, including the strong-field curvature $+\tfrac32
(M/r)^2$; its reciprocal is the proper-time density $\rho_{\rm eff}=\rho_0/
\sqrt{1-2M/r}$ whose leading term is exactly $1\cdot(GM/rc^2)$. \Sderiv{} for the
form.

\subsection{Status and the one external scale}
The amplitude of the $1/r$ well defines the network's effective coupling
$G_{\rm net}=A/w_M$ (well amplitude per unit source weight). Sweeping the source
weight, the density, and the box leaves $G_{\rm net}$ constant (log--log slopes
$-0.000$, $+0.007$, $+0.003$): the coupling is linear and granularity-independent,
ruling out a non-Poisson law. But sweeping the action stiffness gives
$G_{\rm net}\propto1/K$ (slope $-1.000$, reference value $G_{\rm net}=0.981$). The
network therefore derives the \emph{coupling relation}---a linear, Poisson-form law
with exponent $-1$ and amplitude $\propto w_M$---while the numerical \emph{value} of
$G_{\rm net}$ rides on $K$, the action normalisation (the granularity scale).
\Sext{}: the absolute value of $G$ is an input, the form is not. The spatial
dimension $d=3$ entering the shell measure $r^{d-1}$ is likewise geometric input, the
same status the metric had above.

\subsection{What the network adds beyond Poisson}\label{sec:beyondpoisson}

A fair objection must be met head on: have we discovered gravity, or only Poisson's
equation? The minimisation of \emph{any} quadratic gradient action in three
dimensions produces a $1/r$ Green's function---this is standard potential theory. The
question is therefore not whether $1/r$ emerges (it must) but whether the quadratic
action itself is an \emph{output} of the network rather than an \emph{input}.

Three measured facts answer it, each going beyond a generic gradient minimisation.
(a)~The specific quadratic form is not inserted by hand: it is the coarse-grained
Benincasa--Dowker operator on the Poisson sprinkling, whose operator coefficients are
fixed from Poisson isotropy alone (the measured operator ratio $5/9$ to $0.06\%$, with
a cubic lattice exactly blind to the same operator, \S\ref{sec:skyrme}), not chosen.
(b)~The Schwarzschild dilation is reproduced to $0.2\%$ by chain counting
(\S\ref{sec:gravity}) without the metric formula ever entering the generator; this is
more than ``just Poisson'' because it requires the specific causal structure of the
Schwarzschild order, not a flat gradient problem. (c)~The central result
(\S\ref{sec:sourcesgravity}) feeds the soliton's \emph{own} topological energy-density
profile---not a generic point source---into the relaxer, and the exterior exponent
$-0.992$ comes out of that specific defect profile rather than an ad hoc choice of
source shape.

What remains open we state plainly: the gravitational sector established here is
Newtonian, a linear Poisson equation. The non-linear regime---gravitational waves,
back-reaction between sources, and the Einstein equations---lies beyond the present
construction and is not claimed.

A separate, phenomenological line of work asks whether the effective galactic
dynamics of such a vacuum modifies Newtonian gravity at low accelerations, with a
pre-registered, parameter-free prediction for the redshift evolution of the baryonic
Tully--Fisher relation confronted with observations in companion work
\cite{BTFRMNRAS}. That line is independent of the present construction and is not
relied upon here. Within this paper, the gravitational sector stands on the direct
measurement above---the soliton's own relaxed profile sourcing its $1/r$ field in
proportion to its mass---which is what makes the present ``$1/r$ from a network''
more than a restatement of potential theory.

\section{Gauge-group selection: why \texorpdfstring{$\su2$}{SU(2)}}\label{sec:group}

The internal target of Eq.~\eqref{eq:sigma} was provisional. We now ask which group
the requirement of \emph{stable, conserved, point-like matter} forces, and show by
elimination that, among the candidates considered, $\su2$ is the minimal one that
works.

\begin{theorem}[Elimination of candidate internal groups]\label{thm:elim}
Require an internal target space $G$ whose field configurations on $\mathbb{R}^3$
(tending to a constant at spatial infinity, hence classified by $\pi_3$ of the vacuum
manifold) admit a topological charge that is \textup{(i)}~conserved,
\textup{(ii)}~localised, and \textup{(iii)}~carried by a stable finite-energy
configuration. Then, among $\{$single real scalar, finite/discrete group, $U(1)$,
$\su2$, and the higher classical groups$\}$, the minimal group meeting all three is
$\su2$.
\end{theorem}

\begin{proof}[Proof (by elimination)]
Each step is a topological or measured obstruction.
\emph{Scalar.} A single real scalar has a point (or two-point) vacuum manifold; the
relevant homotopy group vanishes, so there is no conserved point-like charge in three
space dimensions. Fails~(i).

\emph{Discrete/finite groups.} A discrete target gives domain-wall (codimension-one)
defects, not localised point-like ones, and carries no continuous current. Fails~(ii).

\emph{$U(1)$.} The vacuum manifold is $S^1$, and the point-like sector in $3$D is
governed by $\pi_2(S^1)=0$: there is no point-like topological charge; $U(1)$ supports
line vortices ($\pi_1(S^1)=\mathbb{Z}$), not particles. We verified the obstruction
directly on the network: every $U(1)$ class function is exactly blind to a putative
vortex core because $\cos2\pi n=1$, so the candidate point-defect carries no
gauge-invariant marker and diffuses away under the dynamics. Fails~(iii) in the
point-like sector.

\emph{$\su2$.} The group manifold is $S^3$ and $\pi_3(S^3)=\mathbb{Z}$: there is a
conserved, integer-valued, localised topological charge---the baryon number $B$---and
\S\ref{sec:soliton} exhibits a stable finite-energy configuration carrying it. $\su2$
meets~(i)--(iii).

\emph{Higher groups.} For $SU(N\ge2)$ one again has $\pi_3=\mathbb{Z}$, so larger
groups are not excluded by topology; they are excluded by minimality, $\su2$ being the
smallest non-Abelian compact group with $\pi_3=\mathbb{Z}$ (Bott periodicity
guarantees no smaller candidate appears in the relevant degree).
\end{proof}

\emph{Status.} \Sderiv{} in the contained sense: the elimination shows that,
\emph{among the candidates considered}, only $\su2$ admits a stable point-like defect
and $U(1)$ provably cannot. It is a minimality statement, not a uniqueness proof over
all conceivable targets; exotic target spaces outside this family were not examined.

\section{The Skyrme term as a necessary external input}\label{sec:skyrme}

The $\su2$ hedgehog of Theorem~\ref{thm:elim} is, by Derrick's theorem
\cite{Derrick1964}, unstable to collapse under the quadratic (Dirichlet) action
alone: a stabilising higher-derivative term is required, the quartic ``Skyrme'' term
\cite{Skyrme1962}. The question that decides the honesty of the whole matter sector is
whether the network \emph{produces} that term with the strength needed to stabilise
the soliton, or whether it must be added. We prove the latter, as a positive
statement.

Write the left-invariant currents $L_i=U^\dagger\partial_i U\in\mathfrak{su}(2)$. The
quadratic (Dirichlet) and quartic (Skyrme) densities are built from the \emph{same}
currents,
\begin{equation}
S=-\tfrac12\,\mathrm{Tr}(L_iL_i),\qquad
K=-\tfrac1{16}\,\mathrm{Tr}\big([L_i,L_j][L_i,L_j]\big).
\label{eq:S2K}
\end{equation}

\begin{theorem}[Pointwise domination of the quartic by the quadratic]\label{thm:KleS}
At every point of space $K(x)\le\tfrac{2}{3}\,S(x)$, and the hedgehog configuration
saturates the bound.
\end{theorem}

\begin{proof}[Proof (Cauchy--Schwarz on the commutator)]
For Lie-algebra elements regarded as vectors under the Killing inner product, the
squared norm of the commutator obeys
\begin{equation}
\big|[L_i,L_j]\big|^2\;\le\;|L_i|^2\,|L_j|^2-(L_i\!\cdot\!L_j)^2\;\le\;|L_i|^2\,|L_j|^2,
\label{eq:cs}
\end{equation}
the first step being the identity $|a\times b|^2=|a|^2|b|^2-(a\cdot b)^2$ applied to
the structure-constant cross product of $\mathfrak{su}(2)\cong\mathbb{R}^3$, the
second dropping the non-negative $(L_i\!\cdot\!L_j)^2$. Summing Eq.~\eqref{eq:cs} over
the antisymmetric pairs $i<j$ and comparing with Eq.~\eqref{eq:S2K}---the quartic a
sum over pairs of the left-hand side, the quadratic a sum over single indices of
$|L_i|^2$---and carrying the index combinatorics of three spatial directions (three
ordered pairs against three single indices) yields
\begin{equation}
K(x)\;\le\;\tfrac{2}{3}\,S(x).
\label{eq:KleS}
\end{equation}
For the hedgehog $U=\exp\!\big(i f(r)\,\hat r\!\cdot\!\vec\tau\big)$ the radial and
angular currents are mutually orthogonal, so $(L_i\!\cdot\!L_j)^2=0$, the first
inequality is an equality, and the bound is saturated.
\end{proof}

\begin{proposition}[Spontaneous dominance is impossible]\label{prop:nodom}
Stabilisation against Derrick collapse requires the quartic to dominate the energy
balance at the core. By Theorem~\ref{thm:KleS} the quartic density can never exceed
$\tfrac23$ of the quadratic density the minimal action already contains; integrated
over space the quartic contribution is bounded above by a fixed fraction of the
quadratic one and cannot become dominant. Hence the stabilising dominance cannot
emerge from the minimal action; it must be supplied with an independent coefficient.
\end{proposition}

\emph{What the network does and does not supply.} The \emph{form} of the Skyrme
operator does emerge from Poisson coarse-graining: a measured operator ratio of $5/9$
to $0.06\%$, with a cubic lattice exactly blind to the same operator, fixes the Skyrme
length $\lambda_{\rm Sk}=a/\sqrt{120}$ from Poisson isotropy alone (\Sderiv{} for the
form). But the \emph{dominance}---the core-cost strength that actually stabilises the
soliton---is bounded away from dominance by Theorem~\ref{thm:KleS} and
Proposition~\ref{prop:nodom}, and is therefore the one declared \Sext{} ingredient of
the matter sector. The bound $K\le\tfrac23 S$ was tested against $10^6$ adversarially
generated $\su2$ configurations chosen to try to violate it; none did. We flag this
prominently: the Skyrme stabiliser is external by a theorem, not by a failed search.

\section{The soliton: mass and spin-statistics}\label{sec:soliton}

With the group forced (\S\ref{sec:group}) and the stabiliser declared external
(\S\ref{sec:skyrme}), we exhibit the matter particle and measure its quantum numbers.

\subsection{Existence, charge, mass}
\emph{Method.} The $\su2$ hedgehog is carried by unit quaternions, so no complex
literal enters and the guard passes verbatim; we relax the field to a stationary
configuration under the quadratic plus externally weighted quartic energy.
\emph{Result.} A stable soliton exists with baryon number $B=1$ and virial balance
$E_2=E_4$ at the minimum, with mass $M\approx146$--$207$ in lattice units across the
configurations tested. \Sderiv{} for existence, stability, and $B=1$.

\subsection{Spin one-half and fermionic statistics}
\emph{Method.} We quantise the soliton's collective rotational coordinate following
Adkins--Nappi--Witten \cite{AdkinsNappiWitten1983} and measure the exchange phase
directly on the lattice. \emph{Result, by a triple check.}
(i)~The rotational spectrum is $E_j\propto j(j+1)$, the rigid-rotor law.
(ii)~The Finkelstein--Rubinstein constraint \cite{FinkelsteinRubinstein1968} on
$\pi_3(S^3)=\mathbb{Z}$ configurations selects half-integer $j$, giving $j=\tfrac12$
as the ground state. (iii)~A $2\pi$ rotation of the field maps the configuration to
minus itself, $\psi\to-\psi$, measured to $5\times10^{-16}$. The exchange statistics
are then measured independently:
\begin{equation}
\begin{split}
[\text{swap two solitons}]&=[\text{rotate one by }2\pi]\\
&=1\in\pi_1(\mathrm{SO}(3))=\mathbb{Z}_2,
\end{split}
\label{eq:spinstat}
\end{equation}
with the general law $\varepsilon(n)=(n-1)\bmod 2$ verified for $n=1,2,3$ solitons.
The winding-$2$ framed-transfer correction $\varepsilon(2)=1$ entering the exchange
reading is uniform across distinct winding-$2$ swap fields (measured on three
independent calibrators of known class; A5, \texttt{RESIDUOS\_A5/}), so it is a
property of the swap topology rather than of one field.
The Skyrmion is therefore a fermion of spin $\tfrac12$, with the spin-statistics
connection obtained rather than assumed. \Sderiv{}.

The identification Skyrmion~$\leftrightarrow$~baryon is \Sident{}, not \Sderiv{}: the
topological charge, spin, and statistics match a nucleon's, but the mass is in lattice
units, not GeV, and the absolute lattice scale is external. We therefore say the
Skyrmion ``has the quantum numbers of a baryon,'' not ``is a baryon of mass $X$.'' We
note one death honestly: the rotational band follows the Casimir law
$m^2\propto J(J+1)$ ($R^2=1.0000000$), i.e.\ a rigid rotor, and \emph{not} the linear
Regge law of a hadronic string.

\section{Quantitative baryon phenomenology}\label{sec:baryon}

Section~\ref{sec:soliton} measured the soliton's topological charge, spin, and
statistics; its quantum numbers match a nucleon's, but the correspondence was left
\Sident{}---the mass in lattice units, no continuum observable computed. We now
sharpen it to \Sderiv{} for the \emph{dimensionless} baryon structure, by quantising
the collective rotation quantitatively and comparing with the standard Skyrme-model
phenomenology of Adkins, Nappi and Witten (ANW) \cite{AdkinsNappiWitten1983,Witten1983}.
The construction introduces exactly one dimensionless coupling (the Skyrme parameter
$e$) and one external scale ($f_\pi$, the analogue of $G$); the programme's universal
pattern then predicts that pure numbers are derived, one number calibrates $e$, and the
absolute scale stays external.

\subsection{The rotor tower and its degeneracies}
Projecting the rigid-rotor spectrum onto the Finkelstein--Rubinstein-allowed
half-integer sector gives a tower whose grand-spin levels carry degeneracies
$(2j+1)^2=[4,16,36]$ for $j=\tfrac12,\tfrac32,\tfrac52$---the nucleon doublet, the
$\Delta$ quartet, and the next level---with spin locked to isospin ($I=J=j$) at each
level. The low-lying spectrum follows the rigid-rotor law $E_\ell\propto\ell(\ell+2)$
to $R^2=0.9999$, and the parameter-free level ratio is
\begin{equation}
\frac{E_{5/2}-E_{1/2}}{E_{3/2}-E_{1/2}}=2.639
\qquad(\text{rigid rotor }8/3=2.667,\ 1.1\%).
\label{eq:rotorratio}
\end{equation}
We certify the rotor law on this accurate low-$\ell$ spectrum and on the pure ratio,
not on a high-$\ell$ fit: the $j(j+1)$ regression \emph{including} the $j=\tfrac52$
level is only $R^2=0.98$, because that high transfer level carries discreteness and
Monte-Carlo spread ($35\%$). The degeneracies and Eq.~\eqref{eq:rotorratio} are the
robust, parameter-free content.

\subsection{The full inertia and one calibration}
The earlier qualitative quantisation (\S\ref{sec:soliton}) used only the
$\sigma$-model zero-mode inertia; the quartic Skyrme term stiffens the rotor and was
omitted there. Measuring both contributions from the lattice profile, the Skyrme part
is $43\%$ of the total, so the full rotational inertia is $1.76\times$ the
$\sigma$-only value---a correction every quantitative number inherits (the rotor
energies the $\sigma$-only treatment implied were $76\%$ too large). With the full
inertia, a \emph{single} calibration---the experimental $N$--$\Delta$ splitting---fixes
the one dimensionless coupling,
\begin{equation}
e=5.39\qquad(\text{ANW }5.45,\ 1.0\%),
\label{eq:ecalib}
\end{equation}
reproducing the dimensionless splitting $(M_\Delta-M_N)/M_N=0.312$ (experiment
$0.314$). No second number is tuned.

\subsection{Parameter-free predictions}
With $e$ fixed, the remaining observables are parameter-free
(Table~\ref{tab:baryon}). The magnetic-moment ratio $\mu_p/\mu_n=-1.51$ lands
\emph{between} ANW ($-1.43$) and experiment ($-1.46$); the isoscalar charge radius is
$0.65$~fm; and $g_A=0.56$ reproduces the famously low minimal-Skyrme value (ANW
$0.61$)---a known limitation of the action, reported for completeness, not a free
parameter.

\begin{table}[t]
\caption{Quantitative baryon phenomenology of the $\su2$ soliton (BQ). The
degeneracies, the rotor ratio, and---after the single $N$--$\Delta$ calibration that
fixes $e$---the magnetic-moment ratio, radius, and $g_A$ are measured from the
lattice soliton's own profile and compared with Adkins--Nappi--Witten and experiment.
Only $e$ is calibrated; everything below the rule is a parameter-free prediction.
$f_\pi$ (the absolute GeV scale) is external and is not listed as a prediction.}
\label{tab:baryon}
\begin{ruledtabular}
\begin{tabular}{lccc}
Observable & This work & ANW & Experiment \\
\colrule
degeneracies $(2j+1)^2$ & $[4,16,36]$ & $[4,16,36]$ & $[4,16,36]$ \\
ratio (rotor $8/3$) & $2.639$ & $2.667$ & --- \\
$e$ (calibration) & $5.39$ & $5.45$ & --- \\
\colrule
$(M_\Delta-M_N)/M_N$ & $0.312$ & --- & $0.314$ \\
$\mu_p/\mu_n$ & $-1.51$ & $-1.43$ & $-1.46$ \\
$r_{\rm iso}$ (fm) & $0.65$ & $0.59$ & $0.81$ \\
$g_A$ & $0.56$ & $0.61$ & $1.27$ \\
\end{tabular}
\end{ruledtabular}
\end{table}

\emph{Status, in the same breath.} What is \Sderiv{} is the \emph{dimensionless}
baryon phenomenology: the multiplet degeneracies, the rigid-rotor ratio, the
magnetic-moment ratio, the radius, and $g_A$---measured pure numbers reproducing ANW to
$\le10\%$ from the lattice soliton's own profile. What stays \Sext{} is the overall GeV
scale: $f_\pi$ is the one external input (the calibrated $f_\pi\approx126$~MeV is the
$F_\pi=2f_\pi$ convention; the dimensionless $e$, which controls every prediction, is
convention-robust). The electromagnetic current decomposition (isoscalar baryon
current, isovector third iso-Noether current) is the standard ANW one, imported as the
Skyrme-term form was imported in \S\ref{sec:skyrme}---the \emph{profile} is the
lattice's, and reproducing ANW's numbers from it validates the integrands. No absolute
mass in GeV is claimed.

\section{The soliton sources gravity}\label{sec:sourcesgravity}

This is the result that unites the two halves of the paper. Sections~\ref{sec:gravity}
and~\ref{sec:soliton} are not two independent successes that happen to coexist on the
same substrate; they are one measurement read twice, because the object whose mass we
measured in \S\ref{sec:soliton} is precisely a localised source for the field of
\S\ref{sec:gravity}.

The logic is a composition of measured laws. The soliton of \S\ref{sec:soliton} is a
conserved, finite-energy, spatially localised lump of integrated weight equal to its
mass $M$. The gravitational sector of \S\ref{sec:gravity} responds to \emph{any} such
localised weight $w$ with a field whose exterior form is $1/r$ for every source shape
(the shape-independent harmonic exterior, $\nabla^2\theta=J$) and whose amplitude is
strictly linear in $w$: the effective coupling $G_{\rm net}=A/w$ was measured constant
in the source weight (log--log slope $-0.000$, coefficient of variation zero), in the
density, and in the box. Identifying the source weight with the soliton mass, $w=M$,
the Skyrmion therefore sources
\begin{equation}
\boxed{\;\theta(r)=\frac{G_{\rm net}\,M}{r}\;}
\label{eq:sourcesgravity}
\end{equation}
the \emph{same} $1/r$ field established in \S\ref{sec:gravity}, with an amplitude set
by the \emph{same} mass measured in \S\ref{sec:soliton}, and growing linearly with it
(two solitons source twice the amplitude, by the exact superposition of
\S\ref{sec:gravity}). Matter and its gravitation are thus a single derived relation:
the topological charge that makes the particle stable is, through its energy density,
the source weight that sets the depth of the gravitational well it sits in.

\subsection{Direct measurement: the soliton's own profile as the source}
The argument just given composes two separately measured laws. We have also performed
the direct measurement it anticipates, feeding the soliton's \emph{own} relaxed
energy-density profile $\varepsilon(r)=\mathcal{E}_2(r)+\mathcal{E}_4(r)$---not a
generic localised weight---as the literal source $s_i$ into the same static BD relaxer,
Eq.~\eqref{eq:BDenergy}. The genuine numerical risk is specific: the soliton's
concentrated, compact profile (unlike the smooth extended shapes of
\S\ref{sec:gravity}) could corrupt the Poisson exterior through grid discreteness, the
conservation offset, and core compactness. It does not. Across the soliton family
($M=175$ to $218$ in lattice units, spanned by varying the external Skyrme weight), the
exterior exponent is $-0.992$ (target $-1$), the well amplitude per unit mass
$G_{\rm net}=A/M=0.9307$ is constant to five digits (coefficient of variation
$5\times10^{-5}$), and it equals the $G_{\rm net}$ of a top-hat source of identical mass
to $0.02\%$. A validation gate first confirmed that the exact relaxer reproduces the
Metropolis BD minimum of \S\ref{sec:gravity}. Equation~\eqref{eq:sourcesgravity} is
therefore a single direct measurement, not an inference: the object whose mass we
measured in \S\ref{sec:soliton} sources the field of \S\ref{sec:gravity} with precisely
that mass.

\emph{Status, stated in the same breath.} The \emph{form}---$\theta\propto M/r$, with
amplitude linear in the mass---is \Sderiv{}, now by direct measurement: the
shape-independent Poisson exterior, the linearity of $G_{\rm net}$ in the source weight,
and the soliton's own profile-as-source all measured. The \emph{absolute}
proportionality constant is \Sext{}: the coupling is $G_{\rm net}\propto1/K$ with $K$
the external action normalisation (\S\ref{sec:gravity}), so
Eq.~\eqref{eq:sourcesgravity} fixes the gravitational \emph{response per unit mass} only
up to that one granularity scale, and the soliton mass itself is in lattice units
(\S\ref{sec:soliton}). One caveat we state plainly: the near-exact $A\propto M$ and the
top-hat agreement are, in part, Gauss's law---the harmonic exterior depends only on the
total enclosed charge---which is precisely why the shape-independence over four extended
profiles held in \S\ref{sec:gravity}; the direct measurement confirms it holds for the
soliton too \emph{and} that the compact profile does not corrupt the exponent. The
radial reduction used here is exact for the embedded $B=1$ hedgehog. As an intermediate
check we have since sourced the soliton energy density on a full \emph{three-dimensional
Cartesian} Poisson grid (no imposed radial symmetry) and solved $-K\nabla^2\theta={}$source
directly: the exterior exponent is $-0.991$ and $G_{\rm net}=A/M$ is constant to
$\mathrm{CV}=0\%$ across four soliton masses ($M=175$--$218$), confirming
$\theta=G_{\rm net}M/r$ is not an artefact of radial symmetry (A5,
\texttt{RESIDUOS\_A5/}). The remaining step---an actual irregular causal sprinkling into
a $3$D BD relaxer---meets the same non-locality obstruction documented for the gauge
sector and is left as an optional larger campaign, the radial and $3$D-Cartesian
compositions being the faithful objects for the present claim.

\emph{A falsifiable prediction.} One non-trivial consequence of this section is the
exact proportionality $G_{\rm net}=A/w_M$, measured constant across the source weight,
the density, and the box size (log--log slopes $-0.000$, $+0.007$, $+0.003$;
\S\ref{sec:gravity}). This rules out a non-Poisson law and establishes that the
gravitational coupling is linear in the source weight with no anomalous dimension---a
result that is \emph{not} guaranteed by the quadratic action alone. It is a concrete,
testable prediction of the construction: any modification that introduced a
non-linearity in the source weight (an amplitude scaling as $w_M^{1+\delta}$ with
$\delta\neq0$) would falsify this sector. The measured $\delta=0$ to three decimals is
thus a parameter-free output the framework stakes itself on.

\section{What is, and is not, established}\label{sec:claim}

\emph{Established by measurement.}
(a)~A localised source in the causal-density field relaxes, under the BD action with
no imposed profile, to a Newtonian $1/r$ potential; the law is Poisson and linear for
every source shape, and a chain-counting measurement reproduces the Schwarzschild
dilation to $0.2\%$ (\S\ref{sec:gravity}).
(b)~Stable, conserved, point-like matter forces the internal group to $\su2$ by
elimination, $U(1)$ being excluded by $\pi_2(S^1)=0$ both by theorem and on the
network (\S\ref{sec:group}).
(c)~The Skyrme stabiliser cannot dominate spontaneously ($K\le\tfrac23 S$, $10^6$
adversarial configurations) and is an external input by theorem (\S\ref{sec:skyrme}).
(d)~The $\su2$ soliton exists with $B=1$, has virial mass $M\approx146$--$207$ in
lattice units, and is a measured spin-$\tfrac12$ fermion (\S\ref{sec:soliton}).
(e)~Quantising its collective rotation with the full ($\sigma$+Skyrme) inertia, a
single calibration (the $N$--$\Delta$ splitting) fixes the one coupling to $e=5.39$
($1\%$ from ANW), after which the magnetic-moment ratio $\mu_p/\mu_n=-1.51$ (experiment
$-1.46$), the isoscalar radius $0.65$~fm and the multiplet degeneracies $[4,16,36]$ are
parameter-free: the dimensionless baryon structure is derived (\S\ref{sec:baryon}).
(f)~Feeding the soliton's \emph{own} energy-density profile as the literal source into
the relaxer of~(a) yields $\theta=G_{\rm net}M/r$ directly---exterior exponent $-0.992$,
$G_{\rm net}$ constant to five digits across the soliton mass, identical to a top-hat of
equal mass to $0.02\%$ (\S\ref{sec:sourcesgravity}).

\emph{Not established; external or identified.} The Skyrmion--baryon correspondence is
\Sderiv{} for the \emph{dimensionless} structure (degeneracies, ratios,
magnetic-moment ratio; \S\ref{sec:baryon}) but the absolute mass scale remains \Sext{}:
$f_\pi$, and hence every mass in GeV, is an external input, exactly as $G$ and the
lattice scale are. The electromagnetic currents are the standard ANW operators,
imported; $g_A$ inherits the $\sim10$--$30\%$ accuracy of the minimal Skyrme action. The
absolute value of Newton's constant ($G_{\rm net}\propto
1/K$) and the lattice scale are \Sext{}; the spatial dimension $d=3$ entering the
shell measure is geometric input. The gravitational sector here is the linear
(Newtonian/weak-field) one; the full $\sqrt{1-2M/r}$ strong field requires the
nonlinear completion and is not claimed. The wave excitation of the orientation field
is treated separately and is a pair of relativistic internal \emph{scalars}
\cite{GoldstonePRD}, not a gauge vector; no electromagnetic or photonic result is
claimed anywhere in this paper. Finally, this paper is deliberately bounded to the
$\su2$ matter-and-gravity sector: $\su3$ colour and confinement, galactic-scale
dynamics, and dark matter are out of scope and are the subjects of separate companion
papers.

\section{Conclusions}\label{sec:conclusions}

Equipping a causal set with a single internal orientation field, and analysing it
under a guard that forbids inserting relativity, we have established by measurement
that the substrate produces both halves of a matter--gravity pair from one
construction. On the gravitational side, a localised source relaxes to a Newtonian
$1/r$ potential obeying a linear Poisson law for every source shape, with the
Schwarzschild dilation recovered to $0.2\%$ by chain counting and the coupling form
derived (its absolute value external). On the matter side, the requirement of stable
point-like topological charge forces $\su2$, the soliton's stabiliser is external by a
Cauchy--Schwarz theorem, and the resulting $\su2$ Skyrmion is a $B=1$,
spin-$\tfrac12$ fermion with measured exchange statistics that is, quantitatively, a
baryon: with the full ($\sigma$+Skyrme) inertia a single calibration fixes the one
coupling, and the magnetic-moment ratio, radius, and multiplet degeneracies then follow
parameter-free, matching Adkins--Nappi--Witten to $\le10\%$ (only the GeV scale
external). The two halves meet in a single \emph{direct} measurement: the soliton's own
energy-density profile, fed as the literal source into the network's gravitational
relaxer, produces $\theta=G_{\rm net}M/r$ with exterior exponent $-0.99$ and amplitude
equal to its own mass. Within its declared external scales, a pure causal substrate with
one internal field yields stable matter that gravitates---matter and gravity as one
measurement, not two coincidences.

\begin{acknowledgments}
This work uses only the author's own simulations; the anti-circularity guard, the
validation gates, and the pre-registered decision criteria are part of the build.
\end{acknowledgments}

\bibliographystyle{apsrev4-2}
\bibliography{TEIC_CONSOLIDATED}

\end{document}
